% Part: counterfactuals
% Chapter: minimal-change-semantics
% Section: antecedent-strengthening

\documentclass[../../../include/open-logic-section]{subfiles}

\begin{document}

\olfileid{cnt}{min}{agg}

\olsection{పూర్వపక్షాన్ని బలపరచడం}

``పూర్వపక్షాన్ని బలపరచడం'' అంటే
$!A \lif !C \Entails (!A \land !B) \lif !C$
అనే నిగమనం. ఇది భౌతిక సోపాధికానికి
చెల్లుతుంది; ప్రతివాస్తవాలకు మాత్రం
చెల్లదు. నేను ఈ అగ్గిపుల్లను గీసి
ఉంటే అది మండి ఉండేదనేది సత్యం
అనుకుందాం. (అంటే అగ్గిపుల్లలో గానీ,
అగ్గిపెట్టె రాపిడి ఉపరితలంలో గానీ
లోపం లేదు; నేను పుల్లను విరగ్గొట్టను;
వగైరా.) కానీ నేను ఈ అగ్గిపుల్లను
బాహ్య అంతరిక్షంలో గీసి ఉంటే అది
మండి ఉండేదనేది సత్యం కాదు.
కాబట్టి కింది నిగమనం చెల్లదు:
\begin{quote}
  అగ్గిపుల్లను గీసి ఉంటే అది
  మండి ఉండేది.

  కాబట్టి అగ్గిపుల్లను బాహ్య
  అంతరిక్షంలో గీసి ఉంటే అది
  మండి ఉండేది.
\end{quote}

లూయిస్--స్టాల్నేకర్ సోపాధిక
వివరణ దీనిని వివరిస్తుంది:
నేను అగ్గిపుల్లను బాహ్య
అంతరిక్షంలో గీసే అత్యంత సమీప
లోకం, దాన్ని సాధారణంగా గీసే
అత్యంత సమీప లోకంతో పోలిస్తే
వాస్తవ లోకానికి చాలా దూరం.
రెండో లోకంలో అగ్గిపుల్ల
మండుతుందనేది సత్యమే అయినా,
మొదటి లోకంలో మండదు.
అదే సరైన ఫలితం.
\sourcecorrection{OLTECNTANT-001}{మూల గద్యంలోని మూడు చోట్ల అగ్గిపుల్లను ‘గీయడం’ బదులు ‘మండించడం’ అని చెప్పడంతో, అగ్గిపుల్ల మండుతుందా అనే ఉత్తరపక్షంతో పూర్వపక్షం కలిసిపోయింది. ఉదాహరణ, ఉటంకించిన నిగమనాలకు అనుగుణంగా మూడింటినీ ‘గీయడం’గా సరిచేశాం.}

\begin{ex}
  గోళ అర్థవిచారంలో ఈ నిగమనం
  చెల్లదు; అంటే $p
  \cif r \Entails/ (p \land q) \cif r$.
  $W = \{w, w_1, w_2\}$,
  $O_w = \{\{w\}, \{w,
  w_1\}, \{w, w_1, w_2\}\}$,
  $V(p) = \{w_1, w_2\}$,
  $V(q) = \{w_2\}$,
  $V(r) = \{w_1\}$
  గల నమూనా $\mModel{M}
  = \tuple{W, O, V}$ను
  పరిశీలించండి. $p$ సత్యమయ్యే
  లోకాన్ని కలిగిన గోళం
  $S = \{w, w_1\}$ ఉంది;
  దానిలోని అన్ని లోకాల వద్ద
  $p \lif r$ సత్యం. కాబట్టి
  $\mSat{M}{p
    \cif r}[w]$. అలాగే
  $(p \land q)$ సత్యమయ్యే
  లోకాన్ని కలిగిన గోళం
  $S' = \{w, w_1, w_2\}$
  ఉంది. కానీ
  $\mSat/{M}{(p \land q) \lif r}[w_2]$;
  అందువల్ల
  $\mSat/{M}{(p \land q) \cif r}[w]$
  (\olref{fig:antecedent-strengthening} చూడండి).

\begin{figure}
\begin{center}
\begin{tikzpicture}[layerwidth=1.5,scale=.6]\tiny
  \spheresystem{3}
  \propositionintersect{-10}{3}{30}{5}
  \begin{scope}
    \clip plot[smooth,tension=1] coordinates
          {(0,4.5) (30:.95) (4.5,-3.5)} -- (4.5,4.5) ;
    \spherefill{2}
  \end{scope}
  \draw[proposition] plot[dashed,smooth,tension=1] coordinates
          {(0,4.5) (30:.95) (4.5,-3.5)} ;
  \proposition{30}{2}{30}{5}
  \path (0,0) node[world] {$w$};
  \spherepos{30}{2}{node[world] {$w_1$}}
  \spherepos{-10}{3}{node[world] {$w_2$}}
  \spherepos{-10}{4}{node {$q$}}
  \spherepos{30}{4}{node {$r$}}
  \draw (1,4.5) node[below] {$p$};
\end{tikzpicture}
\caption{పూర్వపక్ష బలపరచడానికి ప్రతిదృష్టాంతం}
\ollabel{fig:antecedent-strengthening}
\end{center}
\end{figure}
\end{ex}

\end{document}
